\chapter{Kubernetes}
\label{cap:kubernetes}

\begin{inbreve}
I container rendono un'applicazione portabile, ma da soli non dicono che cosa fare quando un container muore, dove
eseguirlo e come farlo comunicare con gli altri. Lo fa una piattaforma di \emph{container orchestration}, e Kubernetes
è la più diffusa. Il capitolo ne spiega i principi di design, gli oggetti con cui si descrive un'applicazione (Pod,
Deployment, Service, Ingress) e i componenti del cluster che li realizzano. Si chiude con i casi in cui Kubernetes non
serve. Presuppone i capitoli~\ref{cap:container} e~\ref{cap:microservizi}.
\end{inbreve}

\section{Perché un orchestratore di container}
\label{sec:k8s-perche}

Si richiamano i vantaggi dei container (capitolo~\ref{cap:container}).
\begin{itemize}
  \item Isolano l'ambiente di un'applicazione con un meccanismo leggero.
  \item Un'immagine si esegue in modo affidabile su qualunque macchina: la stessa immagine va dallo sviluppo alla
    distribuzione (\emph{portabilità}).
  \item Più carichi di lavoro (\ingl{workload}) possono stare sulla stessa macchina fisica, e la memoria e la CPU
    sono usate di più.
\end{itemize}

Quando un'applicazione va in produzione restano però delle domande.
\begin{itemize}
  \item Che cosa succede se il container muore?
  \item Che cosa succede se si guasta la macchina su cui gira il container?
  \item Se più container devono comunicare tra loro, come si realizza la rete tra i container?
  \item Se l'ambiente di produzione è fatto di più macchine, come si decide su quale eseguire un container?
\end{itemize}

A queste domande risponde la \termine{container orchestration}. L'architetto del sistema definisce quali container
eseguire. Poi passa il controllo alla piattaforma di orchestrazione, che si occupa di realizzare e mantenere quella
configurazione.

\begin{definizione}[Kubernetes]
\termine{Kubernetes} (abbreviato \termine{K8s}: la K, otto lettere, la s) è una piattaforma di container
orchestration. Gestisce l'intero ciclo di vita dei singoli container, avviando e spegnendo risorse secondo necessità.
\end{definizione}

Kubernetes risponde alle domande una per una.
\begin{itemize}
  \item Se un container si spegne in modo inatteso, Kubernetes reagisce avviandone un altro al suo posto.
  \item Offre un meccanismo con cui le applicazioni comunicano tra loro, anche mentre i singoli container sottostanti
    vengono creati e distrutti (\rif{sec:k8s-service}).
  \item Dato un insieme di workload da eseguire e un insieme di macchine in un \termine{cluster}, esamina ogni
    container e decide la macchina migliore su cui pianificarlo (\ingl{schedule}).
\end{itemize}
Un cluster è un insieme di macchine collegate in rete che Kubernetes gestisce come un'unica risorsa.

\subsection{Desired state management\aggiunta}
\label{sec:k8s-stato}

L'idea centrale si chiama \termine{desired state management}, gestione dello stato desiderato. Chi distribuisce
un'applicazione non dà a Kubernetes una sequenza di comandi. Scrive in un file YAML lo stato che vuole: quali Pod
eseguire (un Pod è un gruppo di container, \rif{sec:k8s-pod}), con quali immagini e in quante copie (\ingl{replicas}).
I servizi del cluster ricevono il file attraverso un'API e distribuiscono i Pod sulle macchine \emph{worker}, i
\emph{container host}. Su ogni worker un processo, il \emph{kubelet}, avvia e controlla i container
(\rif{sec:k8s-worker}).

\begin{esempio}[Che cosa succede se si perde un worker]
Il file \texttt{App1.yaml} chiede 3 copie del Pod~1 (con due container) e 2 copie del Pod~2 (con un container). Il
cluster ha tre worker e distribuisce le 5 copie tra di loro. Un worker si guasta e le copie che ospitava spariscono.
Lo stato reale non coincide più con quello desiderato: per esempio restano 2 copie del Pod~1 invece di 3.
Kubernetes se ne accorge e avvia le copie mancanti sui worker rimasti. Nessuno deve intervenire a mano.
\end{esempio}

\section{Principi di design}
\label{sec:k8s-principi}

\subsection{Declarativeness}

Si definisce soltanto lo \textbf{stato desiderato} del sistema. Kubernetes:
\begin{itemize}
  \item rileva quando lo stato reale del sistema non corrisponde alle attese;
  \item interviene per correggere il problema.
\end{itemize}
Il sistema diventa così \termine{self-healing}, capace di ripararsi da solo.

Lo stato desiderato è definito da una collezione di \emph{oggetti} (\rif{sec:k8s-oggetti}).
\begin{itemize}
  \item Ogni oggetto ha una \termine{specification}, in cui si scrive lo stato desiderato, e uno \termine{status}, che
    riflette lo stato corrente dell'oggetto.
  \item Kubernetes interroga di continuo (\ingl{polling}) ogni oggetto per verificare che lo status sia uguale alla
    specification.
  \item Se un oggetto non risponde, Kubernetes ne avvia una nuova versione che lo sostituisce.
  \item Se lo status di un oggetto si è allontanato dalla specification, Kubernetes dà i comandi necessari per
    riportarlo allo stato desiderato.
\end{itemize}

\subsection{Distribution}

Kubernetes offre un'interfaccia \emph{unificata} per interagire con un cluster di macchine
(\figrif{fig:k8s-distribuzione}). Non bisogna preoccuparsi di comunicare con ogni macchina singolarmente.

\begin{figure}[ht]
\centering
\begin{tikzpicture}
\foreach \x/\c/\n in {3.2/blu/c1, 4.6/verde/c2, 6.0/rosso/c3}
  \node[draw=\c, fill=\c!15, rounded corners=2pt, font=\footnotesize, minimum width=1.2cm] (\n) at (\x,2.0) {Container};
\node[draw, fill=grigiochiaro, minimum width=8.0cm, minimum height=0.8cm] (k) at (4.6,0.6) {Kubernetes};
\foreach \x/\n in {1.6/n1, 3.6/n2, 5.6/n3, 7.6/n4}
  \node[draw, fill=grigiochiaro, minimum width=1.8cm, minimum height=0.8cm] (\n) at (\x,-0.9) {Node};
\foreach \n in {c1,c2,c3} \draw[->] (\n) -- (\n |- k.north);
\foreach \n in {n1,n2,n3,n4} \draw[<->] (\n |- k.south) -- (\n);
\node[font=\footnotesize, align=right, anchor=east] at (0.3,2.0) {workload\\individuali};
\node[font=\footnotesize, align=right, anchor=east] at (0.3,0.6) {interfaccia\\unificata};
\node[font=\footnotesize, align=right, anchor=east] at (0.3,-0.9) {sistema\\distribuito};
\end{tikzpicture}
\caption{Distribution. I workload (i container in alto) si consegnano a Kubernetes, che offre un'unica interfaccia e
comunica con i nodi del cluster. Chi distribuisce l'applicazione non parla con i singoli nodi.}
\label{fig:k8s-distribuzione}
\end{figure}

\subsection{Decoupling}

Ogni container va sviluppato pensando a un'unica responsabilità (\ingl{single concern}). È l'idea dell'architettura
a microservizi (capitolo~\ref{cap:microservizi}). Kubernetes sostiene in modo naturale i \textbf{servizi
disaccoppiati}, che si possono scalare e aggiornare in modo indipendente.

\subsection{Immutable infrastructure}
\label{sec:k8s-immutabile}

Per ottenere il massimo da container e orchestrazione si distribuisce una \termine{immutable infrastructure},
un'infrastruttura che non si modifica mai sul posto.
\begin{itemize}
  \item Non si entra in un container su una macchina per aggiornare una libreria. Si costruisce una nuova immagine con
    la modifica, si distribuisce la nuova versione e si termina la vecchia.
  \item Durante il ciclo di vita del progetto (sviluppo, test, produzione) si usa \textbf{la stessa immagine}. Si
    cambiano solo configurazioni esterne all'immagine, per esempio montando un file di configurazione.
  \item I container sono pensati come \emph{effimeri}: pronti a essere sostituiti in qualsiasi momento da un'altra
    istanza.
  \item Tornare a uno stato precedente (\ingl{roll back}), per esempio dopo un errore, è facile: basta aggiornare la
    configurazione perché usi un'immagine più vecchia.
\end{itemize}

\section{Gli oggetti di Kubernetes}
\label{sec:k8s-oggetti}

Gli oggetti di Kubernetes si definiscono in file chiamati \termine{manifest}, scritti in YAML o in JSON. Un manifest
ha quattro parti (\tabrif{tab:k8s-manifest}).

\begin{table}[ht]
\centering
\small
\caption{Le quattro parti di un manifest e la domanda a cui risponde ciascuna.}
\label{tab:k8s-manifest}
\begin{tabularx}{0.9\linewidth}{>{\ttfamily}lX}
\toprule
\normalfont Parte & Domanda \\
\midrule
apiVersion & Qual è il gruppo dell'API dell'oggetto e qual è la sua versione? \\
kind & Che cosa si sta creando: un Pod, un Service o altro? \\
metadata & Quali dati identificano in modo univoco l'oggetto (per esempio il nome)? \\
spec & Quale deve essere lo stato dell'oggetto? Quale immagine usa? Quali altre configurazioni? \\
\bottomrule
\end{tabularx}
\end{table}

\subsection{Pod}
\label{sec:k8s-pod}

\begin{definizione}[Pod]
Un \termine{Pod} è l'unità che Kubernetes esegue. È formato da:
\begin{itemize}
  \item uno o più container strettamente collegati;
  \item uno strato di rete condiviso;
  \item volumi del filesystem condivisi.
\end{itemize}
\end{definizione}

Ogni Pod riceve un indirizzo IP proprio (\figrif{fig:k8s-pod}). I container dello stesso Pod condividono quell'indirizzo
e i volumi.

\begin{figure}[ht]
\centering
\begin{tikzpicture}
\draw[fill=grigiochiaro] (0,0) rectangle (4.6,2.6);
\node[font=\footnotesize\bfseries, anchor=north west] at (0.1,2.55) {Pod};
\node[font=\footnotesize, anchor=north east] at (4.5,2.55) {IP address};
\node[draw=rosso, fill=rosso!20, rounded corners=3pt, minimum width=1.6cm, minimum height=0.75cm, font=\footnotesize] at (1.2,1.4) {Container};
\node[draw=verde, fill=verde!20, rounded corners=3pt, minimum width=1.6cm, minimum height=0.75cm, font=\footnotesize] at (3.4,1.4) {Container};
\node[draw=blu, fill=blu!20, minimum width=4.0cm, minimum height=0.45cm, font=\footnotesize] at (2.3,0.45) {Volume};
\end{tikzpicture}
\caption{Un Pod con due container. I due container condividono l'indirizzo IP del Pod (in alto a destra) e il volume
(in basso).}
\label{fig:k8s-pod}
\end{figure}

\subsection{Deployment}
\label{sec:k8s-deployment}

\begin{definizione}[Deployment]
Un oggetto \termine{Deployment} comprende una collezione di Pod definiti da un \emph{template} e un \termine{replica
count}: il numero $n$ di copie del template che si vogliono eseguire. Il cluster cerca sempre di avere $n$ Pod
disponibili.
\end{definizione}

Per esempio, se un Deployment ha replica count 10 e 3 Pod si guastano, Kubernetes pianifica altri 3 Pod su un'altra
macchina del cluster (\figrif{fig:k8s-deployment}).

\begin{figure}[ht]
\centering
\begin{tikzpicture}
\draw[dashed] (-0.4,-0.6) rectangle (9.0,2.4);
\node[font=\footnotesize, anchor=north west] at (-0.35,2.35) {Deployment};
\node[font=\footnotesize, rotate=90] at (-0.1,0.9) {Template};
\foreach \x in {0.4, 5.6} {
  \draw[fill=grigiochiaro] (\x,0) rectangle (\x+3.0,1.8);
  \node[font=\scriptsize\bfseries, anchor=north west] at (\x+0.05,1.78) {Pod};
  \node[font=\scriptsize, anchor=north east] at (\x+2.95,1.78) {IP address};
  \node[draw=rosso, fill=rosso!20, rounded corners=2pt, minimum width=1.1cm, font=\scriptsize] at (\x+0.8,0.95) {Container};
  \node[draw=verde, fill=verde!20, rounded corners=2pt, minimum width=1.1cm, font=\scriptsize] at (\x+2.2,0.95) {Container};
  \node[draw=blu, fill=blu!20, minimum width=2.6cm, minimum height=0.3cm, font=\scriptsize] at (\x+1.5,0.3) {Volume};
}
\node at (4.5,0.9) {$\cdots$};
\node[font=\footnotesize] at (4.5,-0.3) {$n$ repliche};
\end{tikzpicture}
\caption{Un Deployment. Il template descrive un Pod; il Deployment ne mantiene in esecuzione $n$ copie identiche
(repliche), ciascuna con il proprio indirizzo IP.}
\label{fig:k8s-deployment}
\end{figure}

\begin{esempio}[Deployment di un modello di machine learning]
Il manifest seguente esegue 10 istanze di un container che serve un modello di machine learning attraverso
un'interfaccia REST.
\begin{verbatim}
apiVersion: apps/v1
kind: Deployment
metadata:
  name: ml-model-serving
  labels:
    app: ml-model
spec:
  replicas: 10
  selector:
    matchLabels:
      app: ml-model
  template:
    metadata:
      labels:
        app: ml-model
    spec:
      containers:
      - name: ml-rest-server
        image: ml-serving:1.0
        ports:
        - containerPort: 80
\end{verbatim}
Le parti della \texttt{spec} rispondono a quattro domande.
\begin{description}[leftmargin=0pt, style=unboxed]
  \item[\texttt{replicas}.] Quanti Pod devono essere in esecuzione? Dieci.
  \item[\texttt{selector}.] Come si trovano i Pod che appartengono a questo Deployment? Sono quelli con l'etichetta
    \texttt{app: ml-model}.
  \item[\texttt{template}.] Come deve essere fatto un Pod? La sua \texttt{metadata} aggiunge ai Pod l'etichetta
    \texttt{app: ml-model}, così il Deployment trova i Pod da gestire.
  \item[\texttt{containers}.] Quali container girano nel Pod? Uno solo, \texttt{ml-rest-server}, dall'immagine
    \texttt{ml-serving:1.0}, in ascolto sulla porta 80.
\end{description}
\end{esempio}

\subsection{Service}
\label{sec:k8s-service}

Ogni Pod riceve un indirizzo IP unico, con cui si può comunicare con esso. L'insieme dei Pod di un Deployment però
può cambiare in qualunque momento: un Pod si guasta, uno nuovo lo sostituisce con un altro indirizzo. \emph{Come si
comunica con i Pod se il loro insieme cambia di continuo?}

\begin{definizione}[Service]
Un \termine{Service} offre un punto di accesso stabile (\ingl{endpoint}) che dirige il traffico verso i Pod
desiderati. L'endpoint resta lo stesso anche quando i Pod sottostanti cambiano per aggiornamenti, scaling o guasti.
\end{definizione}

Un Service sa a quali Pod mandare il traffico grazie alle \termine{label}: coppie chiave-valore definite nei metadati
dei Pod (\figrif{fig:k8s-service}).

\begin{figure}[ht]
\centering
\begin{tikzpicture}
\node[draw, fill=grigiochiaro, minimum width=3.2cm, minimum height=0.6cm, font=\footnotesize] (s) at (5.0,3.0) {Service};
\draw[->] (5.0,4.0) -- node[right, font=\footnotesize] {traffico} (s.north);
\draw[dashed] (-0.2,-0.9) rectangle (10.2,1.9);
\node[font=\scriptsize, anchor=north west] at (-0.15,1.85) {Deployment};
\foreach \x/\i in {0.2/1, 2.8/2, 5.4/3, 8.0/4} {
  \ifnum\i=2
    \draw[grigio!50, fill=grigiochiaro!40] (\x,-0.4) rectangle (\x+2.0,1.2);
    \node[draw=rosso!40, fill=rosso!8, text=grigio!60, rounded corners=2pt, font=\tiny, minimum width=0.75cm] at (\x+0.5,0.55) {Cont.};
    \node[draw=verde!40, fill=verde!8, text=grigio!60, rounded corners=2pt, font=\tiny, minimum width=0.75cm] at (\x+1.5,0.55) {Cont.};
    \node[font=\scriptsize, text=grigio, anchor=north west] at (\x+0.02,1.18) {Pod};
  \else
    \draw[fill=grigiochiaro] (\x,-0.4) rectangle (\x+2.0,1.2);
    \node[draw=rosso, fill=rosso!20, rounded corners=2pt, font=\tiny, minimum width=0.75cm] at (\x+0.5,0.55) {Cont.};
    \node[draw=verde, fill=verde!20, rounded corners=2pt, font=\tiny, minimum width=0.75cm] at (\x+1.5,0.55) {Cont.};
    \node[font=\scriptsize, anchor=north west] at (\x+0.02,1.18) {Pod};
    \draw[->] (s.south) -- (\x+1.0,1.2);
  \fi
  \node[font=\scriptsize, anchor=north] at (\x+1.0,-0.05) {\texttt{app: ml-model}};
}
\node[font=\footnotesize, text=grigio, align=center, anchor=north] at (3.8,-0.95) {Pod guasto:\\non riceve più traffico};
\node[font=\footnotesize, align=center, anchor=north] at (9.0,-0.95) {Pod nuovo, con la stessa label:\\riceve traffico};
\end{tikzpicture}
\caption{Un Service davanti a un Deployment. Il Service manda il traffico ai Pod con la label \texttt{app: ml-model}.
Il secondo Pod si è guastato (in grigio) e non riceve più traffico. Il Deployment avvia un Pod nuovo (a destra) con la
stessa label, e il Service comincia a mandargli traffico. Chi usa il Service non si accorge del cambio.}
\label{fig:k8s-service}
\end{figure}

\begin{esempio}[Service per il Deployment precedente]
\begin{verbatim}
apiVersion: v1
kind: Service
metadata:
  name: ml-model-svc
  labels:
    app: ml-model
spec:
  type: ClusterIP
  selector:
    app: ml-model
  ports:
  - protocol: TCP
    port: 80
\end{verbatim}
\begin{description}[leftmargin=0pt, style=unboxed]
  \item[\texttt{type}.] Come si vuole esporre l'endpoint? \texttt{ClusterIP} lo rende raggiungibile all'interno del
    cluster.
  \item[\texttt{selector}.] Come si trovano i Pod verso cui dirigere il traffico? Sono quelli con la label
    \texttt{app: ml-model}, cioè i Pod del Deployment dell'esempio precedente.
  \item[\texttt{ports}.] Come parlano i client con il Service? Con TCP sulla porta 80.
\end{description}
\end{esempio}

\subsection{Ingress}
\label{sec:k8s-ingress}

Un Service espone un'applicazione dietro un endpoint stabile, ma \textbf{solo al traffico interno al cluster}. Per
esporla al traffico esterno serve un oggetto \termine{Ingress}. Con un Ingress si sceglie quali Service rendere
pubblicamente disponibili (\figrif{fig:k8s-ingress}). È il modo con cui in Kubernetes si aggiunge un API gateway a
un'applicazione (\rif{sec:smell-microfreshener}).

\begin{figure}[ht]
\centering
\begin{tikzpicture}
\node[font=\footnotesize] (ext) at (5.0,4.4) {traffico esterno};
\draw (-0.3,-1.6) rectangle (10.3,3.7);
\node[font=\footnotesize\bfseries, anchor=north west] at (-0.25,3.65) {Kubernetes cluster};
\node[draw, fill=grigiochiaro, minimum width=3.0cm, minimum height=0.55cm, font=\footnotesize] (in) at (5.0,2.9) {Ingress};
\node[draw, fill=grigiochiaro, minimum width=2.8cm, minimum height=0.55cm, font=\footnotesize] (s1) at (2.4,1.7) {Service};
\node[draw, fill=grigiochiaro, minimum width=2.8cm, minimum height=0.55cm, font=\footnotesize] (s2) at (7.6,1.7) {Service};
\draw[->] (ext) -- (in);
\draw[->] (in.south) |- (5.0,2.3) -| (s2.north);
\draw[<->] (s1) -- (s2);
\foreach \cx/\col in {2.4/rosso, 7.6/blu} {
  \draw[dashed] (\cx-2.4,-0.95) rectangle (\cx+2.4,1.0);
  \foreach \dx in {-1.2, 1.2} {
    \draw[fill=grigiochiaro] (\cx+\dx-1.0,-0.8) rectangle (\cx+\dx+1.0,0.6);
    \node[font=\tiny, anchor=north west] at (\cx+\dx-0.98,0.58) {Pod};
    \node[draw=\col, fill=\col!20, rounded corners=2pt, font=\tiny, minimum width=0.9cm] at (\cx+\dx,-0.15) {Container};
    \draw[->] (\cx,1.42) -- (\cx+\dx,0.6);
  }
}
\node[font=\scriptsize, anchor=north] at (2.4,-0.98) {Deployment};
\node[font=\scriptsize, anchor=north] at (7.6,-0.98) {Deployment};
\end{tikzpicture}
\caption{Ingress. Il traffico esterno entra nel cluster solo attraverso l'Ingress, che lo dirige al Service scelto (a
destra). I due Service comunicano tra loro dentro il cluster; ciascuno distribuisce le richieste ai Pod del proprio
Deployment. Il Service di sinistra non è esposto all'esterno.}
\label{fig:k8s-ingress}
\end{figure}

\begin{esempio}[Ingress]
\begin{verbatim}
apiVersion: networking.k8s.io/v1beta1
kind: Ingress
metadata:
  name: ml-product-ingress
  annotations:
    kubernetes.io/ingress.class: "nginx"
    nginx.ingress.kubernetes.io/rewrite-target: /
spec:
  rules:
  - http:
      paths:
      - path: /app
        backend:
          serviceName: user-interface-svc
          servicePort: 80
\end{verbatim}
\begin{description}[leftmargin=0pt, style=unboxed]
  \item[\texttt{annotations}.] Opzioni di configurazione per l'\emph{Ingress controller}, il componente che realizza
    l'Ingress (qui nginx).
  \item[\texttt{http}, \texttt{paths}.] Come deve accedere al servizio il traffico esterno? Con richieste HTTP al
    percorso \texttt{/app}.
  \item[\texttt{backend}.] A quale Service va diretto il traffico? A \texttt{user-interface-svc}, porta 80.
\end{description}
\end{esempio}

\subsection{Altri oggetti}

Gli oggetti sono molti di più. Per esempio: Job, Volume, Secret, Namespace, ConfigMap, HorizontalPodAutoscaler,
StatefulSet.

\section{Il control plane}
\label{sec:k8s-control-plane}

Il \termine{control plane} è l'insieme dei componenti che controllano il cluster. Un cluster ha due tipi di macchine.
\begin{itemize}
  \item Il \termine{master node}: una macchina (spesso una sola) che contiene la maggior parte dei componenti del
    control plane.
  \item I \termine{worker node}: le macchine che eseguono i workload dell'applicazione.
\end{itemize}
\figrif{fig:k8s-architettura} mostra tutti i componenti, descritti di seguito.

\begin{figure}[ht]
\centering
\begin{tikzpicture}[c/.style={draw, rounded corners=3pt, fill=white, font=\footnotesize, minimum width=2.4cm, minimum height=0.6cm}]
% master
\draw[fill=grigiochiaro] (0,2.4) rectangle (10.4,5.4);
\node[font=\footnotesize\bfseries, anchor=north west] at (0.05,5.35) {Master node};
\node[c] (cm) at (1.8,4.4) {controller-manager};
\node[c] (sc) at (1.8,3.1) {scheduler};
\node[c] (api) at (5.6,3.75) {API server};
\node[draw, cylinder, shape border rotate=90, aspect=0.25, fill=white, minimum width=1.4cm, minimum height=1.0cm, font=\footnotesize] (etcd) at (8.9,3.75) {etcd};
\draw[<->] (cm.east) -- (api.north west);
\draw[<->] (sc.east) -- (api.south west);
\draw[<->] (api) -- (etcd);
% user
\node[font=\footnotesize, align=center] (u) at (5.6,6.5) {utente\\(manifest)};
\draw[->] (u) -- (api.north);
% workers
\foreach \x/\i in {0/1, 3.55/2, 7.1/3} {
  \draw[fill=grigiochiaro] (\x,-0.6) rectangle (\x+3.3,1.35);
  \node[font=\scriptsize\bfseries, anchor=north] at (\x+1.65,-0.65) {Worker node};
  \node[draw, rounded corners=2pt, fill=white, font=\scriptsize] (k\i) at (\x+0.75,0.85) {kubelet};
  \node[draw, rounded corners=2pt, fill=white, font=\scriptsize] (p\i) at (\x+2.4,0.85) {kube-proxy};
  \draw[fill=white] (\x+0.25,-0.45) rectangle (\x+1.55,0.4);
  \node[font=\tiny] at (\x+0.9,-0.02) {Pod};
  \draw[fill=white] (\x+1.75,-0.45) rectangle (\x+3.05,0.4);
  \node[font=\tiny] at (\x+2.4,-0.02) {Pod};
}
\foreach \i in {1,2,3} \draw[<->] (k\i.north) -- (api.south);
\end{tikzpicture}
\caption{Architettura di un cluster Kubernetes. L'utente consegna i manifest all'API server del master node. Nel master
node controller-manager e scheduler lavorano solo attraverso l'API server, che salva lo stato del cluster in etcd. Su
ogni worker node il kubelet chiede all'API server quali Pod eseguire; kube-proxy collega in rete i container dei diversi
nodi.}
\label{fig:k8s-architettura}
\end{figure}

\subsection{Il master node}

\begin{description}[leftmargin=0pt, style=unboxed]
  \item[API server.] L'utente gli consegna la specification di un oggetto nuovo o aggiornato. L'API server valida le
    richieste di aggiornamento e fa da interfaccia unica per ogni domanda sullo stato corrente del cluster.
  \item[etcd.] Lo stato del cluster è conservato in \termine{etcd}, un \ingl{key-value store} distribuito, cioè un
    archivio di coppie chiave-valore replicato su più macchine. Lo stato comprende la configurazione del cluster, le
    specification e gli status degli oggetti, i nodi del cluster, le assegnazioni degli oggetti ai nodi.
  \item[Scheduler.] Lo \termine{scheduler} decide dove eseguire gli oggetti (\figrif{fig:k8s-scheduler}):
    \begin{enumerate}
      \item chiede all'API server quali oggetti non sono ancora stati assegnati a una macchina;
      \item decide a quali macchine assegnarli;
      \item risponde all'API server, che registra l'assegnazione.
    \end{enumerate}
  \item[Controller-manager.] Il \termine{controller-manager} osserva lo stato del cluster attraverso l'API server. Se
    lo stato reale differisce da quello desiderato, apporta modifiche, sempre attraverso l'API server, per portare il
    cluster verso lo stato desiderato. È il componente che realizza il principio di declarativeness.
\end{description}

\begin{figure}[ht]
\centering
\begin{tikzpicture}[c/.style={draw, rounded corners=3pt, fill=grigiochiaro, font=\footnotesize, minimum width=2.4cm, minimum height=0.7cm}]
\node[c] (sc) at (0,0) {scheduler};
\node[c] (api) at (5.0,0) {API server};
\node[draw, cylinder, shape border rotate=90, aspect=0.25, fill=grigiochiaro, minimum width=1.4cm, minimum height=1.0cm, font=\footnotesize] (etcd) at (9.6,0) {etcd};
\draw[->] ([yshift=8pt]sc.east) -- node[above, font=\scriptsize, align=center] {1. ci sono oggetti non\\assegnati a un nodo?} ([yshift=8pt]api.west);
\draw[->] ([yshift=8pt]api.east) -- ([yshift=8pt]etcd.west);
\draw[->] ([yshift=-8pt]etcd.west) -- node[below, font=\scriptsize, align=center] {2. ecco gli oggetti\\non assegnati} ([yshift=-8pt]api.east);
\draw[->] ([yshift=-8pt]api.west) -- ([yshift=-8pt]sc.east);
\draw[->] (sc.south) -- ++(0,-1.0) -| node[pos=0.25, below, font=\scriptsize, align=center] {3. assegna l'oggetto A al nodo X\\e l'oggetto B al nodo Y} (api.south);
\draw[->] (api.north) -- ++(0,0.9) -| node[pos=0.25, above, font=\scriptsize] {4. registrato} (etcd.north);
\end{tikzpicture}
\caption{Lo scheduler al lavoro. (1--2) Chiede all'API server gli oggetti non assegnati, e l'API server li legge da
etcd. (3) Decide i nodi e comunica le assegnazioni all'API server. (4) L'API server le salva in etcd.}
\label{fig:k8s-scheduler}
\end{figure}

\subsection{I worker node}
\label{sec:k8s-worker}

\begin{description}[leftmargin=0pt, style=unboxed]
  \item[kubelet.] Il \termine{kubelet} fa da ``agente'' del nodo.
    \begin{itemize}
      \item Comunica con l'API server per sapere quali workload sono stati assegnati al nodo (``quali Pod devo
        eseguire?'').
      \item Avvia i Pod che eseguono i workload assegnati.
      \item Quando un nodo entra nel cluster, il kubelet ne annuncia l'esistenza all'API server, così lo scheduler può
        assegnargli dei Pod.
    \end{itemize}
  \item[kube-proxy.] Il \termine{kube-proxy} permette ai container di comunicare tra loro attraverso i diversi nodi
    del cluster.
\end{description}

\begin{idea}[Tutto passa dall'API server]
Nessun componente parla direttamente con un altro. Utente, scheduler, controller-manager e kubelet leggono e scrivono
lo stato del cluster solo attraverso l'API server, e solo l'API server accede a etcd.
\end{idea}

\section{Considerazioni finali}
\label{sec:k8s-conclusioni}

Kubernetes \emph{non} conviene:
\begin{itemize}
  \item se il workload può girare su una sola macchina;
  \item se le esigenze di calcolo sono leggere;
  \item se non serve alta disponibilità e si può tollerare che il sistema resti fermo per un po';
  \item se non si prevede di fare molte modifiche ai servizi distribuiti;
  \item se si ha un monolite e non si intende scomporlo in microservizi.
\end{itemize}

Un'alternativa è Docker Swarm, l'orchestratore di Docker (\tabrif{tab:k8s-swarm}).

\begin{table}[ht]
\centering
\small
\caption{Docker Swarm e Kubernetes a confronto.}
\label{tab:k8s-swarm}
\begin{tabularx}{\linewidth}{>{\raggedright\arraybackslash}X>{\raggedright\arraybackslash}X}
\toprule
Docker Swarm & Kubernetes \\
\midrule
più semplice da installare & offre l'auto-scaling \\
curva di apprendimento più dolce & più tollerante ai guasti \\
 & comunità molto grande \\
 & sostenuto dalla Cloud Native Computing Foundation (CNCF) \\
\midrule
preferito dove si privilegiano semplicità e sviluppo veloce & preferito per cluster medi o grandi che eseguono
  applicazioni complesse \\
\bottomrule
\end{tabularx}
\end{table}

\begin{puntichiave}
\begin{itemize}
  \item La container orchestration risponde a ciò che i container da soli non risolvono: container che muoiono,
    macchine guaste, rete tra container, scelta della macchina. Kubernetes (K8s) è una piattaforma di orchestration.
  \item Quattro principi:
    \begin{itemize}
      \item declarativeness: si dichiara lo stato desiderato, K8s confronta specification e status e corregge
        (self-healing);
      \item distribution: un'interfaccia unica per il cluster;
      \item decoupling: un container, una responsabilità;
      \item immutable infrastructure: nuova immagine invece di modifiche sul posto; stessa immagine in sviluppo, test e
        produzione.
    \end{itemize}
  \item Gli oggetti si descrivono in manifest YAML o JSON: apiVersion, kind, metadata, spec.
  \item Pod: container strettamente collegati con rete e volumi condivisi. Deployment: template di Pod più replica
    count. Service: endpoint stabile verso i Pod scelti con le label, solo interno. Ingress: espone Service al
    traffico esterno.
  \item Master node: API server (unica interfaccia), etcd (stato del cluster), scheduler (assegna i Pod ai nodi),
    controller-manager (riporta il cluster allo stato desiderato). Worker node: kubelet e kube-proxy.
  \item Kubernetes non serve per workload piccoli, su una sola macchina, senza bisogno di alta disponibilità o per un
    monolite che resta tale. Docker Swarm è più semplice; Kubernetes è più completo.
\end{itemize}
\end{puntichiave}

\begin{domande}
\begin{enumerate}
  \item Quali domande lasciano aperte i container da soli? Come risponde Kubernetes a ciascuna?
  \item Che cosa significa che Kubernetes è dichiarativo? Che differenza c'è tra specification e status?
  \item Che cos'è un'immutable infrastructure? Perché rende facile il rollback?
  \item Quali sono le quattro parti di un manifest?
  \item Che cos'è un Pod? Che relazione c'è tra Pod e Deployment?
  \item Un Deployment ha replica count 10 e 3 Pod si guastano: che cosa succede?
  \item Perché serve un Service? Come sa a quali Pod mandare il traffico?
  \item Che differenza c'è tra Service e Ingress?
  \item Descrivere i componenti del master node e il ruolo dello scheduler.
  \item Che cosa fanno kubelet e kube-proxy?
  \item In quali casi non conviene usare Kubernetes? Confrontarlo con Docker Swarm.
\end{enumerate}
\end{domande}
